% !TEX root = ../main.tex
%=====================================================================
%  §1  一些背景历史与本讲义的规划
%  §2  后续讲义的规划
%  原文 pp.3--8（PDF pp.4--9）
%=====================================================================

\section{一些背景历史与本讲义的规划}

\begin{keytakeaways}
\item \textbf{物理图像：} 四维流形的分类问题在拓扑层面与光滑层面截然不同：
      弗里德曼定理在拓扑范畴内“完全解决”了单连通四维流形的分类，而唐纳森定理说明
      光滑结构把这一分类彻底打碎——光滑四维流形的世界比拓扑世界狭窄得多、也贫瘠得多。
\item \textbf{数学工具：} 相交形式（幺模对称整系数双线性型）及其分类、怀特海定理、
      塞尔（Serre）对不定型的分类、弗里德曼与唐纳森的定理；以及作为物理进路杠杆的
      低能有效理论（LEET）与度规标度极限。
\item \textbf{核心结论：} 唐纳森多项式是光滑结构的不变量，可用反自对偶瞬子模空间的
      交截理论定义；威滕 1988 年用 $\mathcal{N}=2$ 超对称杨--米尔斯理论给出其路径积分
      解释，而 1994 年塞伯格--威滕的低能有效理论又反过来使唐纳森多项式可用
      塞伯格--威滕不变量显式算出。本讲义的目标正是把这条逻辑链完整展开。
\end{keytakeaways}

让我们先总结四维流形上的一些标准结果。详见教科书 \cite{ref8,ref16,ref17,ref48}。

在本讲义中，$X$ 始终表示一个紧致、连通、可定向、无边的四维流形。

\subsection{基本群}

首先，若 $\pi$ 是任意一个有限展示群，则存在紧致四维流形 $X$ 使得
$\pi_{1}(X)\cong\pi$。\footnote{要证明马尔可夫（Markov）的这个定理，取 $S^{1}\times S^{3}$
的连通和，有限展示中每个生成元对应一个求和项。在连通和中取定一个基点，并取基本群
每个生成元的一个代表元。于是在 $\pi$ 的展示中，每个关系式对应于这些生成元的一个词，
因而对应于连通和中一条嵌入圆周。围绕每条嵌入圆周画一个管状邻域，使该邻域为
$D\times S^{1}$。将这些管状邻域作微小移动使其互不相交。接着作外科手术，把它们换成
形如 $\tilde{D}\times S^{1}$ 的邻域，此时经线在圆盘 $\tilde{D}$ 中变为可缩。最后利用
塞弗特--范坎彭（Seifert--van Kampen）定理即可证明所得流形满足 $\pi_{1}(X)\cong\pi$。
四维是这一切得以发生的最低维数：参见下文附录 D。}
由于群的字问题（word problem）不可判定，这意味着我们无法指望分类所有紧致四维流形。
但我们仍可以指望理解单连通四维流形。

% [Tutor Note: 【易错】这里"四维是能做到这件事的最低维数"一句，靠的是四维流形上
%  连通和/外科手术的自由度：任何有限展示群都能实现为四维流形的 π1，而三维流形
%  的 π1 受到 Heegard 分解的强约束（附录 D 给出证明）。若读者把这句话误读为
%  "四维流形分类更难"，则恰好说反了：正因为在四维 π1 可以是任意有限展示群，
%  分类问题才不可解；反过来，限制到单连通情形，拓扑分类（弗里德曼）反而是完整的。]

\subsection{相交形式}

还有另一个有意思的拓扑不变量，即相交数。可定向紧致四维流形的相交数给出一个不变量：
\begin{equation}
  H^{2}(X;\ZZ)\times H^{2}(X;\ZZ)\;\longrightarrow\;\ZZ .
  \tag{1.1}
\end{equation}
它是 $X$ 的同伦型的不变量。庞加莱对偶说，在自由部分
$H^{2}(X;\ZZ)/\mathrm{Tors}\big(H^{2}(X;\ZZ)\big)$（秩为 $b_{2}(X)$）上它是一个完美配对，
因而对应于一个对称的、整系数的幺模双线性型 $Q_{X}$。

若 $\alpha$ 是一个上同调类，$S(\alpha)$ 为其庞加莱对偶，则这个定向相交数可以写成
好几种形式：
\begin{equation}
  S(\alpha)\cdot S(\beta)
  \;=\;\int_{X}\alpha\beta
  \;=\;\int_{S(\alpha)}\beta .
  \tag{1.2}
\end{equation}
群 $H^{2}(X;\ZZ)/\mathrm{Tors}\big(H^{2}(X;\ZZ)\big)$ 在我们的讨论中将以如下方式出现：
它是德·拉姆上同调类向量空间中的一个格
\begin{equation}
  \bar{H}^{2}(X)\;\subset\;H^{2}_{\mathrm{DR}}(X),
  \tag{1.3}
\end{equation}
定义为具有整周期的那些类构成的集合。当然，这只有在 $X$ 具有微分结构时才有意义。
在那种情形下
\begin{equation}
  \bar{H}^{2}(X)\;\cong\;H^{2}(X;\ZZ)/\mathrm{Tors}\big(H^{2}(X;\ZZ)\big),
  \tag{1.4}
\end{equation}
并且
\begin{equation}
  Q_{X}(\omega_{1},\omega_{2})\;:=\;\int_{X}\omega_{1}\wedge\omega_{2}.
  \tag{1.5}
\end{equation}

% [Tutor Note: 【约定】这里的 (1.3)--(1.5) 隐含了一个关键约定：\bar H^2(X) 是
%  H^2(X;Z)/Tors 在 H^2_DR(X) 中的像，即"整周期类"构成的格。注意它与
%  H^2(X;R) 中庞加莱对偶给出的格是同一个；但在 b_2^+=1 的流形上，\bar H^2(X)
%  中哪些元素是"自对偶类"会依赖于度规诱导的 Hodge 星号，这一点在第 7 节
%  的穿墙讨论中至关重要。另注意 Q_X 的符号约定：此处取 (1.5) 式，若改为
%  -∫ω_1∧ω_2，则下文的 b_2^+ 与 b_2^- 互换，此处与唐纳森原始约定一致。]

相交形式 $Q_{X}$ 的符号差为 $\big(+1^{b^{+}_{2}},\,-1^{b^{-}_{2}}\big)$。也就是说，
若把它看成实向量空间 $H^{2}(X;\ZZ)\otimes\RR$ 上的二次型，则适当选取基之后可化成这一
对角形式。更内蕴地说，$b^{+}_{2}$ 是 $H^{2}(X;\ZZ)$ 中使 $Q_{X}$ 之限制为正定的
极大子格的秩。

当 $X$ 可定向并带有黎曼度规（或仅具共形结构）时，可以定义霍奇对偶
$\ast:\Omega^{2}(X)\to\Omega^{2}(X)$。它满足 $\ast^{2}=+1$，于是可以谈自对偶形式
\begin{equation}
  \ast\omega=\omega
  \tag{1.6}
\end{equation}
与反自对偶形式
\begin{equation}
  \ast\omega=-\omega .
  \tag{1.7}
\end{equation}
霍奇定理允许我们把 $H^{2}(X;\RR)$ 与调和二形式空间 $H^{2}(X)$ 等同起来，而 $\ast$
保持这个空间。于是可以把 $b^{+}_{2}$ 解释为调和自对偶二形式所成向量空间的维数，
把 $b^{-}_{2}$ 解释为调和反自对偶二形式所成向量空间的维数。

\begin{remark*}
\textbf{注记。} 后面我们会与 $\bar{H}^{2}(X)$ 在 $H^{2}(X;\RR)$ 中的各种扭子（torsor）
打交道。例如，给定 $w_{2}(X)$ 的一个整提升 $\bar{w}_{2}(X)$，我们构造
\begin{equation}
  \Gamma_{w}\;:=\;\tfrac{1}{2}\bar{w}_{2}(X)+\bar{H}^{2}(X)\;\subset\;H^{2}(X;\RR).
  \tag{1.8}
\end{equation}
与其他模二类之整提升相伴的类似扭子也将起作用。
\end{remark*}

% [Tutor Note: 【易错】"扭子"（torsor）在这里的意思是：Γ_w 不是一个格（对加法不封闭），
%  而是 \bar H^2(X) 的一个平移副本。这直接导致后文 (5.68) 中通量扇区求和
%  λ∈Γ_ξ 与 (4.75) 中自旋^c 结构 λ∈Γ_w 都是"在一个平移轨道上求和"。
%  计算穿墙相位 (−1)^{(λ−λ_0)·w_2(X)} 时，必须记住 λ−λ_0 ∈ \bar H^2(X)，
%  否则 (−1) 的指数没有良好定义。]

\subsection{怀特海定理}

1949 年，J. H. C. Whitehead（Whitehead）引进了流形的 CW 分解概念，用以分类同伦型。
在 \cite{ref35} 中，米尔诺（Milnor）注意到一个有意思的推论：两个单连通可定向四维流形
$X_{1},X_{2}$ 同伦等价，当且仅当 $Q_{X_{1}}\cong Q_{X_{2}}$。

要证明它，注意单连通四维流形的胞腔分解为
\begin{equation}
  \bigvee_{i}S^{2}_{i}\,\vee\,D^{4},
  \tag{1.9}
\end{equation}
并且由映射 $f:S^{3}\to\bigvee_{i}S^{2}_{i}$ 的同伦类所决定，而这个同伦类可以与相交
矩阵联系起来。\footnote{更详细地说，这样一个映射的同伦类可以这样刻画：对每个 $S^{2}_{i}$
取一个不是楔积基点的点 $p_{i}$，并可使 $p_{i}$ 为 $f$ 的正则值，使得 $f$ 下的原像
是 $S^{3}$ 中的一个纽结。这些纽结有一个环绕矩阵（对角元为自环绕数）。这个环绕矩阵与
$H^{2}(X,\ZZ)$ 上的相交矩阵密切相关。此项事实的解释承 J. Morgan与 G. Segal告知。}

% [Tutor Note: 【省略】米尔诺定理的证明在原文中被压缩成一句话。其真正内容是：
%  单连通的四维 CW 复形由 2-骨架（∪S²_i）与一个 4-胞腔粘合映射 f 决定，
%  而同伦类 [f] ∈ π_3(∨S²_i) 由 Hopf 不变量型的"环绕矩阵"完全刻画，
%  该矩阵正是相交矩阵 Q_X。要点是：dim π_3(∨S²_i) 的生成元与 S²_i 的对
%  一一对应，故 Q_X 决定 [f]，从而决定同伦型。]

\subsection{塞尔的定理}

于是问题归结为整幺模型的分类。塞尔（Serre）对\emph{不定}情形给出了一个漂亮的分类。

\begin{center}
\renewcommand{\arraystretch}{1.25}
\begin{tabular}{|c|c|c|}
\hline
 & 不定（Indefinite） & 定（Definite）\\
\hline
偶（even） & $mE_{8}\oplus nH,\ m\in\ZZ,\ n>0$ & $1,\,2,\,24,\,>10^{7},\dots$\\
\hline
奇（odd）   & $m\langle+1\rangle\oplus n\langle-1\rangle$ & 太多\\
\hline
\end{tabular}
\end{center}

偶定型只存在于维数 $\equiv 0 \pmod 8$ 的情形。上表列出了前几个维数中互不等价的
偶定型的个数。八维中唯一的格是 $E_{8}$ 根格。

$H$ 表示 $\ZZ\oplus\ZZ$ 上由下式给出的偶整型：
\begin{equation}
  Q=\begin{pmatrix}0&1\\1&0\end{pmatrix}.
  \tag{1.10}
\end{equation}

% [Tutor Note: 【易错】表中"偶不定型 = mE_8 ⊕ nH"这一分类要求 n>0（即不定），
%  这是塞尔定理的内容；而偶定型在 dim ≡ 0 mod 8 之外的维数不存在，这一点
%  是 E_8 根格的幺模性与偶性所强制的。读者容易误记成"偶数维才有偶幺模型"，
%  正确说法是"维数被 8 整除（当型正定时）"。另外注意 (1.10) 中 H 的特征多项式
%  λ²−1，其符号差为 0——这解释了后文 11/8 猜想中加 H 不改变 σ 的原因。]

\subsection{弗里德曼定理：同胚型}

迈克尔·弗里德曼（Michael Freedman）在 1982 年的一篇文章中取得了惊人的突破：

\begin{theorem}[弗里德曼]
对每个幺模整型 $Q$，存在单连通紧致可定向拓扑流形 $X$ 使 $Q\cong Q_{X}$。此外：
\begin{enumerate}[label=(\arabic*), leftmargin=2em, itemsep=2pt]
\item 若 $Q$ 为偶型，则这样的 $X$ 在同胚意义下唯一；
\item 若 $Q$ 为奇型，则恰有两个同胚型，而其中至多一个可以是光滑的。
\end{enumerate}
\end{theorem}

% [Tutor Note: 【约定】定理中"拓扑流形"一词不可省：弗里德曼的构造只在 TOP 范畴内
%  成立，所得的 X 一般不允许光滑结构。第 (2) 条中"至多一个可以是光滑的"正是
%  唐纳森定理（§1.6）能够发挥威力的地方：它把这条拓扑上的二选一限制成
%  实际上的唯一。”

要看这个定理有多惊人，注意当 $Q=0$ 时它证明了（四维的、拓扑的）庞加莱猜想。
当 $Q=1$ 时我们有 $X=\CC\PP^{2}$，但必定还存在另一个流形——“假 $\CC\PP^{2}$”——
它与 $\CC\PP^{2}$ 同胚，却不允许任何光滑结构！

\subsection{唐纳森定理：微分同胚型}

几乎与弗里德曼的工作同时（文章均发表于 1983 年），西蒙·唐纳森（Simon Donaldson）
宣布了一些同样惊人的定理。

第一，若 $X$ 允许光滑结构且 $Q_{X}$ 为定，则它必须是对角的：
$m\langle 1\rangle$，其中 $m\in\ZZ$ 非零。

值得注意的推论包括：对应于 $2E_{8}$ 的流形不允许光滑结构。（此前所有已知的判据——
特别是罗赫林（Rokhlin）定理——都容许它可能有光滑结构。）类似地，若 $Q$ 为奇、定、
且非标准（即不是全为 $+1$ 或全为 $-1$ 的对角矩阵），则这两个同胚型都无
光滑结构。

% [Tutor Note: 【省略】唐纳森这一定理的证明在原文中没有给出任何提示。其原始证明
%  走的是"瞬子模空间紧化 + 边界项估计"的路线：在 Q_X 为定时取 k=1 的 ASD 瞬子
%  模空间，它是一个以 X 与若干奇点（对应可约联络）为边界的五维配边；奇点处的
%  局部模型是 CP² 上的锥，于是相交矩阵被强制对角化。这条论证在第 4.5 节
%  脚注 13 中被作者一句话带过，可对照阅读。]

第二，唐纳森引进了他著名的多项式不变量。它们是 $H_{0}(X)\oplus H_{2}(X)$ 上的一列
多项式函数，是 $X$ 之光华结构（光滑结构）的不变量。由这些不变量可得的一个惊人论断是：
不能把一个代数曲面写成一个非平凡的连通和 $X_{1}\#X_{2}$，且满足
$b^{+}_{2}(X_{i})>0$。特别地，它不能写成两个代数曲面的连通和。\footnote{回忆
$b^{+}_{2}(X)$ 是 $H^{2}(X;\RR)$ 的最大正定子空间的秩。}

唐纳森的构造使用秩一规范群 $G=SU(2)$ 与 $G=SO(3)$ 的非阿贝尔规范场论。他利用 $X$ 上
主 $G$ 丛之反自对偶联络模空间上的相交理论来定义这些多项式。现在，方程
\begin{equation}
  F^{+}\;=\;F+\ast F\;=\;0
  \tag{1.11}
\end{equation}
用到了一个黎曼度规。但除了 $b^{+}_{2}=1$ 的流形之外，对度规的依赖会消失。在后一种情形下，
这些多项式在度规空间中是分片常数的，但会穿过墙（wall）跳变。关于它们如何穿墙跳变的
穿墙公式已完全弄清。我们将回到这些重要事实。

% [Tutor Note: 【约定】式 (1.11) 中 F^+ = (F + *F)/2 是自对偶投影，作者写作
%  F + *F = 0（省略了因子 1/2），因为在此式中整体倍率无关紧要。但后文
%  (4.2)、(4.31)、(4.70) 等处凡出现 F^+ 与辅助场 D^+ 比较的地方，因子 1/2
%  就必须小心，否则会差出 2 的幂。建议读者在核对 (4.26) 的虚拟维数与
%  (9.1) 的指数时，统一采用 F^+ := (F+*F)/2 的约定。]

唐纳森的不变量被用来证明关于四维流形光滑结构的一些惊人事实。拓扑四维流形的世界可以
相当狂野。例如，$\RR^{4}$ 上有连续无穷多个互不相同的微分结构。

在这一切进展之后——弗里德曼与唐纳森都获得了 1986 年的菲尔兹奖——人们自然会想，
物理是否扮演了重要角色。毕竟，唐纳森用的正是非阿贝尔规范场论和瞬子，而这两者在现代
粒子理论中起着主要作用。

这正是威滕登场之处。1988 年，他给出了唐纳森多项式的一个量子场论描述
\cite{ref53}。我们将详细说明它（部分地依照 M. Atiyah与 L. Jeffrey 在 \cite{ref3}
中为威滕的文章引入的一个格外漂亮的进路）。一句话：四维流形 $X$ 上的配分函数
$Z_{\mathrm{DW}}$（在作用量中加入某些可观测量的插入之后）是 $X$ 上所有唐纳森多项式的
生成函数。技术上，威滕的量子场论被称为“$G$ 杨--米尔斯理论的 $\mathcal{N}=2$ 超对称
扩充”，其中 $G$ 单且秩为一。

% [Tutor Note: 〔原稿边注〕作者在此处旁注 "Comment on motivation from M. Atiyah [2]?"，
%  意为他想补充一段关于阿蒂亚 \cite{ref2} 中动机的评论，但尚未写出。译文照录
%  正文，不擅自补写。]

威滕的解释是漂亮的——它是拓扑扭曲这一概念、乃至拓扑场论这一更广概念的滥觞——但当时
并不清楚，把唐纳森多项式解释成一个路径积分究竟能在数学上得到什么。问题在于，四维
相互作用量子场论的路径积分在数学界看来是一种神话般的存在；即便对愿意相信其存在的
物理学家而言，有效地进行计算看起来也是力所不及的。

然而，由于理论是拓扑的，路径积分 $Z_{\mathrm{DW}}$ 不依赖于 $X$ 上的度规。但如果我们
“把 $X$ 上的度规放大”，意即作替换
\begin{equation}
  g_{\mu\nu}\;\to\;t\,g_{\mu\nu}
  \tag{1.12}
\end{equation}
并研究极限 $t\to+\infty$，那么所有长度尺度都趋于无穷。在物理中，长度尺度与能量尺度
由不确定性原理相联系：
\begin{equation}
  L\;\sim\;\frac{\hbar}{E},
  \tag{1.13}
\end{equation}
所以把长度标度到无穷，等同于研究“远红外”——理论在那些与真空只有无穷小差异的过程下的
行为。而这恰恰是可以借助所谓低能有效理论（LEET）给出理论的另一种完全不同描述的情形。

% [Tutor Note: 【易错】这一步"放大度规"是全篇的枢纽，但有两处极易含糊：
%  (i) 它依赖理论的拓扑不变性（T_{μν} 为 Q-恰当，见 (4.20)），而不是依赖
%      "路径积分对度规不敏感"这种朴素说法；
%  (ii) 放大度规会同时把所有能标压到零，因此低能有效理论描述的是 u 平面附近
%      的涨落——但"附近"的尺度是由 Λ 与 X 的体积共同决定的。原文随后在第 5 节
%      反复强调的"极限次序问题"（u→±Λ² 与 t→∞ 的次序）就出在这里。]

LEET 用法的一个好例子是 QCD：在短距离（相对于相互作用变强的那个长度尺度）上，
强作用由夸克和胶子描述；而在长距离上，它仅由 $\pi$ 介子描述。短距离理论是规范群为
$SU(3)$ 的非阿贝尔杨--米尔斯理论，极小耦合到 $N_{f}$ 份基础表示之直和的狄拉克旋量；
长距离理论则是靶空间为 $SU(N_{f})$ 的非线性 $\sigma$ 模型。（在自然界中，
$N_{f}=2$ 或 $N_{f}=3$，取决于 LEET 要多精确。）

对我们的目的而言，只要对 $\mathcal{N}=2$ $SU(2)$ SYM 的 LEET 有足够好的理解，
就有望把唐纳森--威滕路径积分改写成可能带来新洞见的新形式。

这正是实际发生的事。1994 年春天，塞伯格与威滕弄清了 $\mathcal{N}=2$ $SU(2)$ SYM 的
LEET \cite{ref47,ref46}。这些信息足以让威滕把唐纳森多项式惊人地重述为一组新的
（更易处理）四维流形不变量——所谓“塞伯格--威滕不变量”\cite{ref55}。他写下的公式
被称为威滕猜想，并为克龙海默与姆罗夫卡此前发现的唐纳森不变量结构定理 \cite{ref21}
补上了关键的一块。

% [Tutor Note: 〔原稿边注〕作者旁注"Introductory remarks took 36 minutes."（引言部分讲了
%  36 分钟），是讲课计时备注，与内容无关。]

% [Tutor Note: 【约定】此处"威滕猜想"在数学文献中即指 (9.1) 式。需要提醒的是：
%  原文把 (9.1) 叫做"key statement of [55]"，而把"Witten's conjecture"这个名字归
%  给数学界用法；两者是同一件事。物理上它是 IR 展开的直接推论，数学上它在
%  b_2^+>1 且单型的情形已被证明（菲汉--莱内斯等），但在一般情形仍是猜想。]

本讲义的目标是：
\begin{enumerate}[leftmargin=2.2em, itemsep=2pt]
\item 解释威滕关于唐纳森多项式的形式场论解释，即用 $\mathcal{N}=2$ $SU(2)$ SYM 来表达；
\item 解释塞伯格与威滕关于 $\mathcal{N}=2$ SYM 低能动力学的物理洞见，如何导出一个
      （在我看来）令人信服的威滕猜想的推导；
\item 解释物理观点如何说明关于唐纳森多项式的许多结果，乃至引出一些数学预言。
\end{enumerate}

关于唐纳森不变量的物理进路，一些重要参考文献是：
\begin{enumerate}[leftmargin=2.2em, itemsep=2pt]
\item Cordes、Moore、Ramgoolam 的 Les Houches 讲义 \cite{ref4}；
\item 拉巴斯蒂达与马里诺的教科书 \cite{ref24}；
\item 两部 IAS 文集：P. Deligne 等，\emph{Quantum Fields and Strings: A Course for
      Mathematicians}，两卷，AMS 1999。特别相关的是第二卷中 E. Witten的讲义：
      Lectures 1--3、8、12，尤其是 Lectures 16--19。
\end{enumerate}

% [Tutor Note: 【建议】第 3 条所列的威滕 IAS 讲义（尤其是 Lectures 16--19）是学习
%  拓扑扭曲与 Mathai--Quillen 形式的最佳原始材料，与本书第 3、4 节对照阅读，
%  可以把"超空间上的 Berezin 积分"与"微分形式上的 Thom 类"这两套语言对应起来。]

\section{后续讲义的规划}

\begin{keytakeaways}
\item \textbf{物理图像：} 本讲义剩下的部分构成一条完整的逻辑链：从拓扑场论的
      形式结构出发，到把唐纳森理论嵌入其中，再到用低能有效理论求值。
\item \textbf{数学工具：} Mathai--Quillen 形式、局域化、等变上同调、瞬子模空间、
      塞伯格--威滕方程、模形式与椭圆曲线。
\item \textbf{核心结论：} 第 3--4 节建立 UV 理论，第 5--9 节完成 UV→IR 映射并导出
      威滕猜想，第 10--12 节给出应用与预言。
\end{keytakeaways}

\begin{enumerate}[leftmargin=2.2em, itemsep=2pt]
\item 上同调拓扑场论的形式结构：路径积分的 Mathai--Quillen 形式与局域化。
\item 唐纳森理论如何嵌入 MQ 框架：拓扑扭曲的 $\mathcal{N}=2$ SYM。
\item SW 解与真空结构：把可观测量从 UV 映射到 IR。
\item 希格斯分支贡献的一般形式。
\item 库仑分支的求值：$u$ 平面积分。
\item 推导唐纳森不变量与 SW 不变量之间的关系。
\item 单型（simple type）。
\item 物理观点的应用：（例如：超共形单型与广义诺特不等式。）
\item 可能的未来方向。
\end{enumerate}

\subsection{致谢}

% [Tutor Note: 【约定】以下致谢按原文逐句译出，人名保留原文拼写，不音译，
%  以免与文献表 [1]--[58] 中的写法不一致。]

这些讲义是在 2012 年 3 月于西蒙斯几何与物理中心举办的“数学与物理研究生讲习班”
（我与 D. Freed、C. Teleman 共同组织）上一系列讲座的基础上大幅扩充和更新而成的。
我感谢 L. Alvarez-Gaumé、K. Fukaya、J. Morgan与 S. Donaldson邀请我在 2017 年 3 月
于 SCGP 作系列讲座，并感谢他们在讲座期间的许多有益评论。我感谢 CERN 理论组在讲义
更新期间给予的招待。我也感谢 D. Brennan、J. Clingempeel、A. Daemi、A. Dey、
E. Diaconescu、A. Khan、Z. Komargodski、C. LeBrun、T. Mainiero、I. Nidaiev、D. Park、
M. Rocek 与 D. Sullivan 在讲座期间提出的有益问题与评论。我还从 2017 年 7 月 17--18 日
汉堡 Pre-String-Math 学校的学生那里得到了一些极好的问题，并感谢 G. Arutyunov、
E. Meir、I. Runkel 与 C. Schweigert 邀请我在该校讲演。本工作受 Rutgers 的 DOE 项目
DE-SC0010008 资助。最后但同样重要的是，我感谢 M. Mariño与 E. Witten在本讲义主题上
与我的出色合作。
